\documentclass[11pt,twoside]{article}
\usepackage{latexsym}
\usepackage{amssymb,amsbsy,amsmath,amsfonts,amssymb,amscd}
\usepackage{dsfont}
\usepackage{graphicx,color}
\usepackage[colorlinks=true]{hyperref}
\usepackage{epsfig,color}
\usepackage{dsfont}
\usepackage{algorithm}
%\usepackage{algorithm2e}
%\usepackage{algoritmic}
%\usepackage{algpseudocode}

\usepackage[amsmath,framed,thmmarks]{ntheorem}

\usepackage{placeins}
\usepackage{subcaption}
\usepackage{amsmath} 
\numberwithin{equation}{section}
\DeclareMathOperator{\sign}{sign}
\usepackage{longtable, array, booktabs}
\newcommand\red[1]{\color{red}#1} 
\usepackage{tabularx}
\usepackage{mathtools}
\setlength{\oddsidemargin}{0mm}
\setlength{\evensidemargin}{0mm}
\setlength{\topmargin}{5mm}
\setlength{\textheight}{20cm}
\setlength{\textwidth}{15cm}
\renewcommand{\baselinestretch}{1.0}
\parindent 10pt

\DeclareMathOperator{\Tr}{Tr}

\DeclareMathOperator*{\argmax}{arg\,max}
\DeclareMathOperator*{\argmin}{arg\,min}

\newcommand*\diff{\mathop{}\!\mathrm{d}}

\theoremstyle{plain}
\newtheorem{theorem}{Theorem}[section]
\newtheorem{lemma}[theorem]{Lemma}
\newtheorem{definition}[theorem]{Definition}
\newtheorem{remark}[theorem]{Remark}
\newtheorem{proposition}[theorem]{Proposition}
\newtheorem{corollary}[theorem]{Corollary}

\theoremstyle{nonumberplain}
\newtheorem{proof}{Proof.-}



\newcommand{\bd}{\mbf{d}}
\newcommand{\bu}{\mbf{u}}
\newcommand{\bv}{{\mathbf{v}}}
\newcommand{\bw}{\mbf{w}}
\newcommand{\bx}{{\mathbf{x}}}
\newcommand{\bX}{\mbf{X}}
\newcommand{\bbX}{\mathbb{X}}
\newcommand{\bY}{\mbf{Y}}
\newcommand{\mH}{\mathcal{H}}

\newcommand{\rhoT}{\rho_T}
\newcommand{\rhoL}{\rho_T^L}
\newcommand{\rhoLM}{\rho_T^{L,M}}
\newcommand{\brhoL}{\bm{\rho}_T^L}
\newcommand{\rhsF}{\mbf{f}}

\newcommand{\brho}{\boldsymbol{\rho}}
% customized commands


\newcommand{\mres}{\mathbin{\vrule height 1.2ex depth 0pt width
0.13ex\vrule height 0.13ex depth 0pt width 0.9ex}}

\definecolor{myblue}{rgb}{0.0352,0.4981,0.6509}
\newcommand{\GER}{\color{myblue}}

\newcommand{\mbf}[1]{\boldsymbol{#1}}
\newcommand{\lil}[1]{\ell_{#1}}
\newcommand{\innerp}[2]{\langle #1,#2 \rangle}
\newcommand{\inp}[1]{\langle{#1}\rangle}
\newcommand{\dbinnerp}[1]{\langle\hspace{-1mm}\langle{#1}\rangle\hspace{-1mm}\rangle}
\newcommand{\abs}[1]{\big| #1 \big|}
\newcommand{\dquo}[1]{``#1"}
\newcommand{\bxm}{\bx^{m}}
\newcommand{\dotbxm}{\dot\bx^m}
\newcommand{\bXm}{\bX^{m}}
\newcommand{\dotbXm}{\dot\bX^{m}}
\newcommand{\Xtrain}{\bX_{\mathrm{tr}}}
\newcommand{\brm}{\bar{\br}^{m}}
\newcommand{\nbrm}{r^{m}}

% define norms
\newcommand{\norm}[1]{\left\| #1 \right\|}
\newcommand{\Tnorm}[1]{\left\| #1 \right\|_{\mathcal{T}}}
\newcommand{\tnorm}[1]{\left\| #1 \right\|_{\mathcal{T}_t}}
\newcommand{\HDnorm}[1]{d_{\text{HD}}{(#1)}}
\newcommand{\ltworho}[1]{\norm{#1}_{L^2(\rho_T)}}
% define some spaces in real
\newcommand{\realR}[1]{\mathbb{R}^{#1}}
\newcommand{\real}{\mathbb{R}}
\newcommand{\compC}[1]{\mathbb{C}^{#1}}
\newcommand{\comp}{\mathbb{C}}
\newcommand{\integ}{\mathbb{Z}}
% for some special variables
% the vector varaible


\newcommand{\br}{\mbf{r}}

\newcommand{\bB}{\mbf{B}}

%\newcommand{\bX}{\mbf{X}}



\newcommand{\balpha}{\mbf{\alpha}}
%
\newcommand{\mA}{\mathcal{A}}
\newcommand{\mC}{\mathcal{C}}
\newcommand{\mE}{\mathcal{E}}
\newcommand{\mF}{\mathcal{F}}
\newcommand{\mG}{\mathcal{G}}
\newcommand{\mL}{\mathcal{L}}
%\newcommand{\mH}{\mathcal{H}}

\newcommand{\mI}{\mathcal{I}}
\newcommand{\mK}{\mathcal{K}}
\newcommand{\mR}{\mathcal{R}}
\newcommand{\mS}{\mathcal{S}}
\newcommand{\mT}{\mathcal{T}}
\newcommand{\mX}{\mathcal{X}}
\newcommand{\mY}{\mathcal{Y}}
%
\newcommand{\bigO}{\mathcal{O}}
%
\newcommand{\R}{\real}
\newcommand{\dimamb}{D}
\renewcommand{\dim}{d}
\newcommand{\numcl}{K}
\newcommand{\idxcl}{k}
\newcommand{\spaceM}{S}
\newcommand{\cl}{C}
\newcommand{\clof}{\mathpzc{k}}
\newcommand{\adj}{\mathcal{A}}
\newcommand{\eadj}{\adj^E}
\newcommand{\aadj}{\adj^A}
\newcommand{\xiadj}{\adj^{\xi}}
% the coupling strength terms
\newcommand{\cst}{\mathcal{\kappa}}
\newcommand{\vcst}{\cst^{\bv}}
\newcommand{\xicst}{\cst^{\xi}}
\newcommand{\bcst}{\bm{\cst}}
\newcommand{\bvcst}{\bm{\cst}^{\bv}}
\newcommand{\bxicst}{\bm{\cst}^{\xi}}
\newcommand{\bL}{\mbf{L}}
\newcommand{\e}{\epsilon}

%

\newcommand{\st}[1]{\textcolor{red}{#1}}
\newenvironment{resultrevision}{\par\begingroup\color{red}}{\par\endgroup}
\newcommand{\XR}[1]{\textcolor{blue}{#1}}



\title{An asymptotic preserving scheme for kinetic equations describing bacterial traveling pulses}
\author{Canan Akkoyunlu\thanks{Department of Mathematics and Computer Science, İstanbul Kültür
University, 34158 Bakırköy, İstanbul, Türkiye (c.kaya@iku.edu.tr)} \and Jose A. Carrillo\thanks{Mathematical Institute, University of Oxford, Woodstock Road, Oxford, OX2 6GG, UK (jose.carrillo@maths.ox.ac.uk)}\and Gissell Estrada-Rodriguez\thanks{Universitat Politecnica de Catalunya, Barcelona, 08034, Spain. Centre de Recerca Matematica, 08193 Bellaterra, Barcelona, Spain (gissell.estrada@upc.edu)} \and  Xinran Ruan\thanks{School of Mathematical Sciences, Capital Normal University, Beijing, China, 100048 (xinran.ruan@cnu.edu.cn). 
%\newline \noindent email: c.kaya@iku.edu.tr, carrillo@maths.ox.ac.uk, gissell.estrada@upc.edu, xinran.ruan@cnu.edu.cn.\newline \noindent 
} 
}
\date{}



\begin{document}

\maketitle
\begin{abstract}
We consider a kinetic model for bacterial chemotaxis that describes the propagation of traveling pulses driven by a nutrient and a self-produced chemoattractant. In the diffusive regime, the model converges to a macroscopic drift–diffusion system with bounded chemotactic fluxes. We construct a fully discrete asymptotic preserving (AP) finite difference scheme based on a micro–macro decomposition of the kinetic equation. The method is stable uniformly with respect to the scaling parameter and, as this parameter tends to zero, reduces to a consistent discretization of the limiting macroscopic equations. We prove positivity and mass conservation for the limiting scheme, and use a positivity-preserving correction for the kinetic solver. Numerical simulations reproduce the experimentally observed traveling pulses at both the kinetic and the macroscopic level, validate the theoretical traveling wave speed, and illustrate the convergence of the kinetic model to its macroscopic limit.
\end{abstract}

\section{Introduction}


Chemotaxis is the directed movement of organisms in response to chemical fluctuations in the environment. This behavior is essential for survival \cite{franca2006navigating} and plays a crucial role in many biological processes in multicellular organisms, such as wound healing, fertilization and inflammation \cite{rickert2000leukocytes},  while unicellular organisms use it as a foraging mechanism \cite{franca2010g,willard2006signaling}. Many bacteria, in particular, navigate toward favorable conditions by sensing gradients of nutrients or chemoattractants, and away from harmful substances.

Although bacterial chemotaxis has been studied since the end of the 19th century, the first precise mathematical model was proposed by Keller and Segel \cite{keller1971model}. They introduced a coupled system of partial differential equations (PDEs) for the cell density $\rho(t,x)$ and the chemical concentration $S(t,x)$:
\begin{subequations}\label{eq:chemicals}
\begin{align}
\frac{\partial \rho}{\partial t} &= \nabla\cdot (D_\rho (\rho, S) \nabla \rho - \chi (\rho, S) \nabla S) + f(\rho, S),
  \label{eq: KS1}\\
\frac{\partial S}{\partial t} &= D_S \Delta S + g (\rho, S) - h(\rho, S)S.
  \label{eq: KS2}
\end{align}
\end{subequations}
Here $D_\rho(\rho,S)$ denotes the bacterial diffusivity, $\chi(\rho,S)$ the chemotactic
sensitivity, and $f$, $g$, and $h$ describe cell growth and death, chemical production,
and degradation, respectively.


We consider the migration of \textit{E.~coli} in response to a nutrient $N$ and a self-produced chemoattractant $S$, motivated by the microchannel experiments in \cite{saragosti2010mathematical}. In these experiments, approximately $10^5$ cells are initially confined near the left end of a $2\,\mathrm{cm}$ channel, while nutrient is supplied from the right. Within about $100\,\mathrm{s}$, the bacteria organize into a dense pulse that propagates toward the nutrient source with nearly constant speed and mass. The pulse exhibits a characteristic asymmetric profile, with a diffuse front and a steep rear, which is maintained during propagation. Since the duration of the experiment is short compared with the cell division cycle, bacterial growth is neglected.





Following the kinetic descriptions of Alt \cite{alt1980biased} and Dunbar and Othmer
\cite{othmer1988models}, we model the bacterial population by the mesoscopic density
$f(t,x,v)$, where $x \in \mathbb{R}$ and $v \in V$ denotes the velocity. Since we are
interested in planar traveling waves, the dynamics are governed by the one-dimensional
kinetic equation
\begin{equation}
    \partial_t f(t,x,v)+{v}\partial_x f(t,x,v)= \int_{V}\Bigl(T(t,x,v,v')f(t,x,v')-  T(t,x,{v}',v)f(t,x,v)\Bigr)\diff v'\ .\label{eq: main kinetic}
\end{equation}
 The turning kernel $T$ describes the change in direction during the tumbling phase of the movement. $T$ depends on the space and time variables through the chemical signals in the medium,  and $T(\cdot, v,v')$ represents the change of velocity from $v'$ to $v$. For simplicity we are going to assume that $T$ only depends on the current velocity with the condition that $\int_VT(t,x,v)\diff v=1$.
 
 Following \cite{saragosti2010mathematical}, we consider the influence of two chemical species: a chemotactic signal $S(t,x)$ released by the bacteria, and a nutriet $N(t,x)$ consumed by the them. We assume that both chemical cues contribute additively to the turning frequency. 
More precisely, we define
\[
\widetilde{\psi}_1(t,x,v)
=
\chi_S \phi\!\left(\frac{DS}{Dt}\Big|_v\right)
+
\chi_N \phi\!\left(\frac{DN}{Dt}\Big|_v\right),
\qquad
\frac{D}{Dt} = \partial_t + v \partial_x,
\]
where $\phi$ is the intracellular response function and $\chi_S$, $\chi_N$ are the
chemotactic sensitivities. Since the velocity average of $\widetilde{\psi}_1$ does not
vanish in general, we introduce the centered perturbation
\begin{equation}\label{eq: centered psi1}
\psi_1(t,x,v)
=
\widetilde{\psi}_1(t,x,v)
-
\frac1{|V|}\int_V \widetilde{\psi}_1(t,x,w)\,dw,
\end{equation}
which satisfies $\int_V \psi_1(t,x,v)\,dv = 0.$
The rescaled turning kernel is then defined by
\[
T_\varepsilon(t,x,v)
=
|V|\,
\frac{1+\varepsilon \widetilde{\psi}_1(t,x,v)}
{\displaystyle \int_V \bigl(1+\varepsilon \widetilde{\psi}_1(t,x,w)\bigr)\,dw},
\]
and admits the first-order expansion
\begin{equation}\label{eq: T expansion}
T_\varepsilon(t,x,v)=1+\varepsilon \psi_1(t,x,v)+O(\varepsilon^2),
\end{equation}
with
$\int_V T_\varepsilon(t,x,v)\,dv = |V|.$







 

The chemoattractant and nutrient concentrations satisfy the reaction--diffusion equations
\begin{subequations}
\label{eq:SN_intro}
\begin{align}
\partial_t S &= D_S \partial_{xx} S - \alpha S + \beta \rho, \\
\partial_t N &= D_N \partial_{xx} N - \gamma \rho N,
\end{align}
\end{subequations}
where
\[
\rho(t,x) = \int_V f(t,x,v)\,dv\ ,
\]
denotes the macroscopic bacterial density.
Here the chemoattractant is produced by the bacteria at rate $\beta$ and naturally degraded at rate $\alpha$, while the nutrient is consumed at rate $\gamma$.

The main objective of this paper is to construct and validate a numerical scheme that accurately captures the dynamics of the kinetic model for all values of the scaling parameter $\varepsilon$, and recovers the corresponding macroscopic drift–diffusion system in the diffusive limit $\varepsilon\to0$. In particular, we investigate whether the traveling pulses observed at the macroscopic level are consistently reproduced by the kinetic description. To this end, we develop an asymptotic preserving (AP) scheme \cite{filbet2010class,jin1999efficient}, which remains stable as $\varepsilon\to0$ and, for fixed discretization parameters, automatically reduces to a consistent discretization of the limiting macroscopic equations. Our approach builds on the micro–macro decomposition introduced in \cite{lemou2008new} (see Section~\ref{sec: micro macro decomp}). Other AP approaches for kinetic chemotaxis models can be found in \cite{CY13,LWZ18,CKLO19} and the references therein.



In the diffusive limit, the kinetic model formally converges to the macroscopic system
\begin{subequations}
\label{basicmodel}
\begin{align}
\partial_t \rho
&= D_\rho \Delta \rho
   - \nabla \cdot \bigl( \rho u_S + \rho u_N \bigr), \\
\partial_t S
&= D_S \Delta S - \alpha S + \beta \rho, \\
\partial_t N
&= D_N \Delta N - \gamma \rho N,
\end{align}
\end{subequations}
where the drift velocities $u_S$ and $u_N$ are derived from the microscopic response function and take the form
\begin{equation}
\label{eq:flux_intro}
u_S
=
\chi_S J_\phi(-\varepsilon \partial_t S, |\nabla S|)
\frac{\nabla S}{|\nabla S|},
\qquad
u_N
=
\chi_N J_\phi(-\varepsilon \partial_t N, |\nabla N|)
\frac{\nabla N}{|\nabla N|}.
\end{equation}
 Here $\phi$ is the so-called response function
 (Section~\ref{sec: maro system}), which
describes how a single bacterium reacts to changes in the concentration of the
chemoattractant $S$ (or of the nutrient $N$) along its trajectory. The function $J_\phi$
encodes, at the macroscopic level, these microscopic features of the individual response.
The parameter $\varepsilon$ is the ratio between the pulse speed and the speed of individual
bacteria; in the experiments $\varepsilon \approx 0.1$, so the pulse travels one order of
magnitude slower than each individual.

In contrast with the classical Keller–Segel fluxes, these drift terms remain uniformly bounded. This provides a more realistic description of chemotactic motion and allows an analytical characterization of traveling pulses \cite{calvez2019chemotactic}: the flux $u_S$ aggregates the cells into a pulse, while $u_N$ drives this pulse toward higher nutrient concentrations. In the stiff-response limit, where $\phi$ is a step function, the chemotactic drift reduces to $\chi\,\mathrm{sign}(\partial_x S)$. For this case, in \cite{calvez2023nonlinear} the authors studied the one-dimensional model without nutrient and with $S$ in quasi-stationary equilibrium, and proved the nonlinear stability of its stationary states, which are exponentially decaying clusters, i.e. zero-speed pulses.












The paper is organized as follows. Section~\ref{sec: maro system} derives the macroscopic drift–diffusion system from the underlying kinetic chemotaxis model and recalls the theoretical traveling wave speed. Section~\ref{sec: micro macro decomp} introduces the micro–macro decomposition used to construct the asymptotic preserving formulation. In Section~\ref{sec:scheme}, a fully discrete AP finite difference scheme is developed, and its asymptotic preserving, positivity, and mass conservation properties are analyzed. Section~\ref{sec: numerics} presents numerical experiments validating the traveling pulse dynamics, the dependence of the wave speed on the chemotactic sensitivities, and the convergence of the kinetic model towards its macroscopic limit as the scaling parameter tends to zero. Finally, Section~\ref{sec: conclusions} summarizes the main results and discusses possible directions for future work.


% Let 
% \begin{equation}\label{TepsFeps}
% %\begin{cases}
%     T_\epsilon(t,x,v)=\frac{\psi_\epsilon(t,x,{v}) }{\displaystyle\int_{V}\psi_\epsilon(t,x,{v})  \diff {v}}\ ,
%  %   \end{cases}
% \end{equation}
% where we assume that the dependence of the turning operator on the chemical concentration $S$ and nutrient $N$ happens as a perturbation of magnitude $\epsilon\ll 1$ in the following way
% \begin{equation}\label{expansion_psi}
%     \psi_\epsilon(t,x,{v})=1+\epsilon\tilde{\psi}_1(t,x,{v})\ ,
% \end{equation}
% %\paragraph{Assumption 1}
% \begin{equation}
% \tilde{\psi}_1(t,x,{v}) = \phi_1[S]({v}) + \phi_2[N]({v}) \ . \label{hypo2}
% \end{equation}
% Here $\phi_1$ and $\phi_2$ are given by
% $$\phi_1[S]({v})=\chi_S\phi\Bigl(\frac{DS}{Dt}\bigg\rvert_{v}\Bigr)\ \ ,\ \ \phi_2[N](v)= \chi_N\phi\Bigl(\frac{DN}{Dt}\bigg\rvert_{v}\Bigr)\ ,$$
% and $\frac{D}{Dt}=\partial_t+{v}\partial_x$ is the total derivative along the direction given by ${v}$.

% Using \eqref{expansion_psi}, the turning kernel $T_\epsilon$ writes
% \begin{align}
%     T_\epsilon(t,x,{v})=& 1+\epsilon\Bigl[\chi_S\phi\Bigl(\frac{DS}{Dt}\bigg\rvert_{v}\Bigr)+ \chi_N\phi\Bigl(\frac{DN}{Dt}\bigg\rvert_{v}\Bigr)\Bigr]  + {\mathcal{O}(\epsilon^2)} \ ,\label{eq: turning operator scaled}
% \end{align}
% {where the error terms contained in $\mathcal{O}(\epsilon^2)$ integrate to $0$ in the velocity variable. We also note that using \eqref{hypo2} we have $\int_V T_\epsilon(t,x,v)\diff v=1$ which describes the conservation of individuals during the velocity reorientation.}



%As mentioned before, we are going to study this biological phenomenon using an approach distinct from classical modelling. In particular, the biased component in the equation of $\rho$ in (\ref{ks_model}), which is going to be derived from a kinetic description of the motion, will be uniformly bounded, whereas in the case of K-S this biased component is the chemotactic sensitivity function $\chi(\rho, S)$, which is proportional to the gradient of the chemical, is generally unbounded when instability due to aggregation takes place. This will provide us with a more precise and rigorous description of the evolution of baterial movement in this specific context, what is to say, it will allow us to describe analytically the traveling wave behaviour it is observed experimentally.






\section{Macroscopic system of equations}\label{sec: maro system}
In this section we derive the macroscopic system that is obtained from \eqref{eq: main kinetic} following \cite{saragosti2010mathematical}. We introduce two relevant parameters: the maximum speed of a single bacterium $c = \max\{|v|: v \in V\}$ and the mean turning frequency $\lambda_0$.
The main unknown is the speed of the traveling pulse, denoted by $\sigma$. We start by rescaling the kinetic model (\ref{eq: main kinetic}) into a non-dimensional form using the non-dimensional variables, denoted with tildes, as
\begin{equation}
    \tilde{t}= \frac{{t}}{\bar{t}}\ , \quad \tilde{x} = \frac{x}{\bar{x}}\ , \quad {\tilde{V} = \frac{V}{c}}\ , \quad \tilde{T}=\frac{T}{\lambda_0}\ ,
    \nonumber
\end{equation}
to model the traveling wave with $\bar{x} = \sigma\bar{t}$. According to the experiments, the pulse
speed $\sigma$ is much lower than the speed of a single bacterium $c$, and this difference is measured by the ratio $\varepsilon = \sigma / c$, where $\varepsilon \ll 1$.



The kinetic equation can be written in the following form, where we  drop the tilde for simplicity of notation 
{
\begin{equation}
    \varepsilon \partial_t f^\varepsilon+{v}\partial_x f^\varepsilon = \frac{\mu}{\varepsilon}\Bigl[\int_V T_\varepsilon(t,x,v')f^\varepsilon(t,x,v')\diff v'-|V|T_\varepsilon(t,x,v)f^\varepsilon(t,x,v)\Bigr]\ , \label{eq: inital model with scaling}
\end{equation}}where $\mu/\varepsilon=\lambda_0 \bar{x}/c$ and with velocity space \( V = [-1,1] \) (one-dimensional case). 
Here 
\begin{equation}\label{eq: T app}
T_\varepsilon(t,x,v)
=
1+\varepsilon\psi_1(t,x,v).
\end{equation} 
is the first-order approximation of \eqref{eq: T expansion}
{and \(\psi_1\) denotes the centered first-order perturbation defined in \eqref{eq: centered psi1}.}
We consider a regular expansion of $f^\varepsilon$ in $\varepsilon$ up to second order
\begin{equation*}
f^\e(t,{x},{v})=f^0(t,{x},{v})+\varepsilon f^1(t,{x},{v})+\varepsilon^2 f^2(t,{x},{v})+\mathcal{O}(\varepsilon^3)\ , 
\end{equation*}
where $\int_V f^\e(t,{x},{v})\diff {v} = \int_V f^0(t,{x},{v})\diff {v} = \rho(t,{x})$, therefore $\int_V f^i(t,{x},{v})\diff {v}=0$, $\forall i\geq 1$.





Substituting this ansatz into \eqref{eq: inital model with scaling} and collecting powers of $\varepsilon$
\iffalse
{\begin{align*}
    \varepsilon^2\partial_t f^0&+\varepsilon{v}\partial_x\big(f^0+\varepsilon f^1  \big) 
    =-\psi_0(v) f^0+\int_V\psi_0(v')f^0(v')\diff v'\\
    &+\varepsilon\int_V\Bigl(\psi_0(v')f^1(v')+\psi_1(t,x,v')f^0(v')-\psi_0(v)f^1(v)-\psi_1(t,x,v)f^0(v) \Bigr)\diff v'+{\mathcal{O}(\varepsilon^3)}
      \ .
\end{align*}}

\noindent Identifying the different equations in powers of $\varepsilon$, we obtain
\fi
\begin{align}
    \varepsilon^0: & \quad f^0(t,{x},{v})=\frac{\rho(t,{x})}{|V|}\ ,\label{eq: varepsilon0}\\
    \varepsilon^1: & \quad 
    v \partial_x f^0
=
\mu \left(
\int_V \psi_1(v') f^0(v')\,dv'
- |V| \psi_1(v) f^0(v)
- |V| f^1(v)
\right)\ ,\label{eq: varepsilon1}\\
    \varepsilon^2: & \quad \partial_tf^0+{v}\partial_x f^1 =\mu\Bigl(\int_V\psi_1(v')f^1(v')\diff v'-|V|\psi_1(v)f^1(v)-|V|f^2(v)\Bigr) \ \label{eq: varepsilon2}.\
\end{align}
Next, after replacing $f^0$ given by  \eqref{eq: varepsilon0}, we multiply \eqref{eq: varepsilon1} by ${v}$ and we integrate again with respect to ${v}$ to obtain
\begin{align}
    \mu|V|\int_V vf^1(t,x,v) \diff v=&-\partial_x\frac{\rho(t,x)}{|V|}\int_V v^2\diff v-\mu\rho(t,x)\chi_S\int_V\phi(\varepsilon\partial_tS+v\partial_xS)v\diff v\nonumber\\ &-\mu\rho(t,x)\chi_N\int_V\phi(\varepsilon\partial_tN+v\partial_xN)v\diff v.\label{eq: mean direction}
\end{align}
We define
\begin{equation}
\label{eq: moment}
\begin{gathered}
D_\rho
=
\frac{1}{\mu |V|^2}\int_V v^2\,\diff v,
\qquad
u_S
=
-\frac{\chi_S}{|V|}
\int_V \phi(\varepsilon\partial_t S+v\partial_x S)\,v\,\diff v,
\\[4pt]
u_N
=
-\frac{\chi_N}{|V|}
\int_V \phi(\varepsilon\partial_t N+v\partial_x N)\,v\,\diff v .
\end{gathered}
\end{equation}
In the limiting case where the internal response function is $\phi(Y)=-\phi_0\sign(Y),$
the flux terms reduce, up to the constant factor determined by
$\phi_0$ and the velocity measure, to \cite{calvez2019chemotactic}
\[
u_S=\chi_S\sign(\partial_x S),
\qquad
u_N=\chi_N\sign(\partial_x N).
\]
Finally, integrating \eqref{eq: varepsilon2} in ${v}$ and using \eqref{eq: moment} we obtain
\[
\partial_t\rho=D_\rho\partial_{xx}\rho-\partial_x(\rho u_S+\rho u_N)\ .
\]
{
The diffusion coefficient is determined by the velocity distribution and
the turning frequency as in \eqref{eq: moment}. 
In particular, when the normalized scaling \(\mu |V|=1\) is used, this
coefficient becomes
\begin{equation}
\label{def:Drho}
  D_\rho
=
\frac1{|V|}
\int_V v^2\,dv .  
\end{equation}
For example, when \( V = [-1,1] \), we have $D_\rho = \frac{1}{3}$.
}
% In particular, since \( V = [-1,1] \), we have $|V|=2$ and $
% \int_{-1}^1 v^2 dv = \frac{2}{3}.
% $
% Hence, $D_\rho=
% \frac{1}{6\mu}$ and if, for instance, \( \mu = 1 \),
% $D_\rho = \frac{1}{6}$.
We refer to \cite{emako2015existence,emako2016traveling} for more details of the derivation.
Together with \eqref{eq:SN_intro} for $S$ and $N$, this system gives a traveling wave solution where the speed is given by the following theorem.

\begin{theorem}[\cite{calvez2019chemotactic}]\label{thmacro}
There exist a speed $\sigma>0$, and two positive values $N_- < N_+$, such that the system \eqref{basicmodel} admits traveling wave solutions ($\rho (z), S(z), N(z)$) for $z=x - \sigma t$ with the following boundary conditions:
\begin{equation}
    \lim_{|z| \to \infty} \rho = \lim_{|z| \to \infty} S = 0, \hspace{0.6cm} \lim_{z \to -\infty} N = N_- < N_+ = \lim_{z \to + \infty} N.
    \nonumber
\end{equation}
Assume in addition that $\partial_z S$ changes sign only once, and that $\partial_z N$ is positive everywhere. Then, the speed $\sigma > 0$ is uniquely determined by the following implicit relation,
\[
\chi_N - \sigma = \chi_S \frac{\sigma}{\sqrt{4 D_S \alpha + \sigma^2}}.
\]
\end{theorem}
The proof in \cite{calvez2019chemotactic} relies on the ansatz that the chemoattractant $S$ has a unique maximum at $z=0$ and that the nutrient $N$ is increasing. This ansatz decouples system \eqref{basicmodel}, so that one can solve explicitly for $\rho$ and then for $S$; validating the ansatz a posteriori yields the implicit relation for $\sigma$ stated in Theorem \ref{thmacro}.










\iffalse
Now we non-dimensionalised the system. We introduce a new set of variables $\tilde{t}=t/\bar{t}$ and $\tilde{x}=x/\bar{x}$ where $\bar{t}$ and $\bar{x}$ are the macroscopic time and space scales. Then we have
\begin{align}
&\partial_t\rho =\tilde{D}_\rho\partial_{xx}\rho-\partial_x (\rho \tilde{u}_S+\rho \tilde{u}_N)\ ,\label{eq: rho nondimension}\\
&\partial_tS =\tilde{D}_S\partial_{xx} S-\tilde{\alpha} S+\tilde{\beta} \rho\ ,\\
&\partial_t N=\tilde{D}_N\partial_{xx} N-\tilde{\gamma} \rho N\ \label{eq: N nondimension},
\end{align}
where $\tilde{u}_S=\frac{\bar{t}\chi_S}{\bar{x}}\sign(\partial_x S)$, $\tilde{u}_N=\frac{\bar{t}\chi_N}{\bar{x}}\sign(\partial_x N)$ and 
\[
\tilde{D}_\rho=\frac{\bar{t}D_\rho}{\bar{x}^2}\ ,\ \ \tilde{D_S}=\frac{\bar{t}D_S}{\bar{x}^2}\ ,\ \ \tilde{\alpha}=\bar{t}\alpha\ ,\ \ \tilde{\beta}=\bar{t}\beta\ ,\ \ \tilde{D}_N=\frac{\bar{t}D_N}{\bar{x}^2}\ ,\ \ \tilde{\gamma}=\gamma \bar{t}\ .
\]
From now on, for simplicity of notation we are going to drop the tilde in the new coefficients.
\fi

%\section{traveling pulses under competing fluxes}

%\textcolor{red}{(To Canan: In this section we should include the traveling wave solution for the kinetic and macroscopic systems. Please, we should write things clearly and not copy and paste results from previous papers. Since everything is already done in \cite{saragosti2010mathematical} for the macroscopic equations and in \cite{calvez2018well} for the kinetic equation we should only include here the main results: traveling wave solution, speed of propagation etc)}

\iffalse
\section{Existence of traveling wave solutions (not sure if this is necessary)}

In this section we summarise the results obtained in \cite{calvez2019chemotactic} for the existence of traveling wave solutions for the kinetic equation \eqref{eq: main kinetic},  and for the macroscopic system \eqref{basicmodel}.

Consider a traveling wave solution $z=x-ct$ where $c$ denotes the speed of the wave. We rewrite \eqref{eq: main kinetic} in the moving frame variable $z$, {\color{red} ($\sigma$ is used for wave speed in other sections)}
\begin{equation}
\begin{aligned}\label{system}
(v-c)\partial_zf(z,v)&=\int T(z,v'-c)f(z,v')d\nu(v')-T(z,v-c)f(z,v)\ ,\\
-c\partial_z S(z)&=D_S\partial_{z}^2S(z)-\alpha S(z)+\beta \rho(z)\ ,\\
-c\partial_zN(z)&=D_N\partial_{z}^2N(z)-\gamma\rho(z) N(z)
\ ,
\end{aligned}
\end{equation}
where $\nu$ is a symmetric probability measure on $\mathbb{R}^d$, where $V=\textnormal{supp}\ \nu$. Throughout this section we assume that $\nu$ is absolutely continuous with respect to the Lebesgue measure. The probability density function belongs to $L^p$ for some $p > 1$, i.e.
\[
\diff \nu (v)=w(v)\diff v\ ,\ \quad w\in L^p\ , \quad p\in(1,\infty]\ .
\]
We assume that the the fresh nutrient is located on the right so the wave moves toward high concentrations of the nutrient, i.e.  $\partial_zN>0$. We expect the chemoattractant concentration to have a maximum at the cell density peak, and therefore we look for a solution where  $\partial_zS$ changes sign. Without loss of generality we assume that this maximal point is at $z=0$.  We summarise these conditions as follows
\begin{equation}
\begin{aligned}
 &S(\pm \infty)=0\ ,\quad N(\pm\infty)=N_\pm\ ,\\
&\partial_z S>0  \quad  \textrm{for}\quad z<0\ , \quad\textnormal{and}\quad\partial_z S<0 \quad \textrm{for}\quad z>0\ ,\\
&\partial_zN>0\ ,\quad\forall z\in\mathbb{R}\ .\label{eq: cond chemicals}
\end{aligned}
\end{equation}
Moreover, the tumbling rate $T$ in \eqref{system} can take only four values depending on the sign of the gradients,
\begin{eqnarray}\label{tumbl}
   T(z,v-c)=1-\chi_S \textrm{sign}((v-c)\partial_z S(z))-\chi_N \textrm{sign}((v-c)\partial_z N(z))\ ,
\end{eqnarray}

\paragraph{Macroscopic traveling pulse}
We present a mathematical description of traveling pulses at the macroscopic scale. For the proof of this theorem we refer the reader to \cite{calvez2019chemotactic}.

\begin{theorem}[\cite{calvez2019chemotactic}]\label{thmacro}
There exist a speed $c>0$, and two positive values $N_- < N_+$, such that the system (\ref{basicmodel}) admits traveling wave solutions ($\rho (x- c t), S(x - c t), N(x - c t)$) with the following boundary conditions:
\begin{equation}
    \lim_{|z| \to \infty} \rho = \lim_{|z| \to \infty} S = 0, \hspace{0.6cm} \lim_{z \to -\infty} N = N_- < N_+ = \lim_{z \to + \infty} N.
    \nonumber
\end{equation}
Assume in addition that $\partial_z S$ changes sign only once, and that $\partial_z N$ is positive everywhere. Then, the speed $c > 0$ is uniquely determined by the following implicit relation,
\[
\chi_N - c = \chi_S \frac{c}{\sqrt{4 D_S \alpha + c^2}}.
\]
\end{theorem}
{\color{red} The wave speed formula does not match the traveling speed from the numerical solution of the kinetic model.}

For the proof of this theorem we assume that the chemoattractant $S$ has a unique maximum at $z = 0$ and that the nutrient $N$ is increasing. Then, this initial ansatz allows us to decouple the system of equations \eqref{basicmodel}, enabling us to solve for $\rho$, and consequently for $S$. Finally, we will validate the ansatz, which will lead to the equation for $c$ as stated in Theorem \ref{thmacro}.

\paragraph{Kinetic traveling pulse} It is natural to expect that some of the qualitative behaviors observed at the macroscopic scale, such as the existence of traveling pulse solutions, should also be present at the mesoscopic level. This motivates the study of the existence of traveling wave solutions directly within the mesoscopic framework, aiming to confirm that such wave-like structures are inherent to the dynamics of bacterial motion, regardless of the scale under consideration.
In the following, we summarise the results in \cite{calvez2019chemotactic} for completeness.

The space variable of the moving pulse at speed $c$ is denoted by $z = x-ct$.
We use the following convention: the superscripts $^+$ and $^-$ denote the position of the velocity $v$ with respect to the speed $c$: for $v>c$ and for $v<c$, respectively. Whereas the subscripts $_+$ and $_-$ denote the relative position with respect to the origin: for $z>0$ and for $z<0$, respectively. 


Similarly, we denote the spatial density of the same group of bacteria $f^-_+(z,v)$ as 
\[
    \rho^-_+(z) = \int_{v<c} f_+(z,v)d\nu(v).
\]
Moreover, the tumbling rate can take only four values $1\pm\chi_S\pm\chi_N$ denoted by $T^+_+, T^-_+, T^+_-, T^-_-$ which depend on weather $v$ crosses $c$ or weather $z$ crosses the origin or not. We illustrate this convention in Figure \ref{tumbling rates} in Appendix \ref{app: T}.




Therefore, $T\in[1-\chi_S-\chi_N, 1+\chi_S+\chi_N]$ and we conclude that the tumbling frequency is higher when bacteria move in the favorable direction and smaller when bacteria move in the unfavorable direction.

\begin{theorem}[\cite{calvez2019chemotactic}] \label{theorem_existence_mesos}
 Assume $(\chi_S, \chi_N) \in (0, 1/2) \times [0, 1/2)$. Under the assumption (A1), and conditions 
\begin{equation} 
    \frac{c_* + \sqrt{(c_*)^2 + 4 \alpha D_S}}{c_* + \sqrt{(c_*)^2 + 4 \alpha D_S} + 2D_S \lambda_-(c_*)} \le \text{explicit constant},
    \label{constraint1}
\end{equation}

\begin{equation}
    \frac{-c^* + \sqrt{(c^*)^2 + 4 \alpha D_S}}{-c^* + \sqrt{(c^*)^2 + 4 \alpha D_S} + 2D_S \lambda_+(c^*)} \le \text{explicit constant},
    \label{constraint2}
\end{equation}

\begin{equation}
    \text{Either } \chi_N \ge \chi_S \hspace{0.5cm} \text{or} \hspace{0.5cm} \frac{\sqrt{\alpha / D_S} + \lambda_+ (0)}{\sqrt{\alpha / D_S} + \lambda_- (0)} \le \text{explicit constant},
    \label{constraint3}
\end{equation}

where $\lambda_-(0)$ and $\lambda_+(0)$
 are the smallest positive roots of, respectively 
 \begin{equation}
     \int \frac{v}{T_-(v) + \lambda v} d\nu (v) = 0, \hspace{0.4cm} \int \frac{v}{T_+(v) - \lambda v} d \nu (v) = 0.
     \nonumber
 \end{equation}
 There exist a non negative c, and positive functions $(f, S, N) \in (L^1 \cap L^\infty (\mathbb{R} \times V)) \times \mathcal{C}^2(\mathbb{R}) \times \mathcal{C}^2 (\mathbb{R})$ solution of the traveling wave problem,
 \begin{equation}
         \left \{
    \begin{aligned}
       (v-c) \partial_z f(z,v) = \int T(z, v'-c)f(z,v') d\nu (v') - T(z, v-c)f(z,v),\\
       -c \partial_z S(z) - D_S \partial_z^2 S(z) + \alpha S(z) = \beta \rho (z),\\
       -c \partial_z N(z) - D_N \partial_z^2 N(z) = - \gamma \rho (z) N(z),
    \end{aligned}
  \right .
  \label{eq_sys_mesoscopic}
\end{equation}
where the tumbling rate $T$ is given by 

\[
    T(z, v-c) = 1 - \chi_S \sign((v-c)\partial_zS(z)) - \chi_N \sign((v-c)\partial_zN(z)) .
\]

\end{theorem}
The details of the proof of this theorem can be found in \cite{calvez2019chemotactic}.

\noindent An important issue discussed in \cite{calvez2018well} is regarding the uniqueness of the traveling wave, where the authors numerically observed that traveling wave exists for several values of $c$, in a parameter regime which violates conditions \eqref{constraint1}-\eqref{constraint3}. This clearly contradicts the results observed at the macroscopic level, where there exists a unique wave speed.  

\textcolor{red}{(this is something we should look into in our AP scheme)}
\fi

{
\section{Micro-macro decomposition}
\label{sec: micro macro decomp}

Throughout this section, we use the scaling $\mu |V|=1$. 
To construct an asymptotic preserving scheme consistent with the diffusive limit, we introduce the (constant) equilibrium distribution 
\[
M(v)=\frac1{|V|},
\qquad \int_V M(v)\,dv=1.
\]
Since the leading-order term is given by \(f^0=\rho/|V|\), we decompose the kinetic density as
\begin{equation}\label{f_expansion}
f^\varepsilon(t,x,v)=\frac{\rho^\varepsilon(t,x)}{|V|}+\varepsilon g^\varepsilon(t,x,v),
\qquad
\int_V g^\varepsilon(t,x,v)\,dv=0.
\end{equation}


We denote by \(\langle \cdot \rangle\) the integration with respect to \(v\) over \(V\), namely
\[
\langle h\rangle = \int_V h(v)\,dv,
\]
and define the projection operator
$
\Pi h = \langle h\rangle M.
$ Then \((I-\Pi)g^\varepsilon=g^\varepsilon\), where $I$ is the identity.

Integrating \eqref{eq: inital model with scaling} with respect to \(v\), and using the conservation property of the $T_\varepsilon$ operator together with
$
\langle vM\rangle = \frac1{|V|}\int_V v\,dv = 0,
$
we obtain the macroscopic equation
\[
\partial_t \rho^\varepsilon + \partial_x \langle v g^\varepsilon\rangle = 0.
\]

% Next, we define 
% \[
% \mathcal{Q}_\varepsilon(f)
% =
% \int_V T_\varepsilon(t,x,v')f(v')\,dv'
% -
% |V|T_\varepsilon(t,x,v)f(v).
% \]
% % The normalized turning kernel admits the first-order expansion
% % \[
% % T_\varepsilon(t,x,v)
% % =
% % 1+\varepsilon\psi_1(t,x,v)+O(\varepsilon^2),
% % \qquad
% % \langle \psi_1\rangle=0.
% % \]
% In the following construction, we use the first-order approximation of \eqref{eq: T expansion} given by
% \[
% T_\varepsilon^{\rm app}(t,x,v)
% =
% 1+\varepsilon\psi_1(t,x,v).
% \]
% It preserves the normalization
% $
% \langle T_\varepsilon^{\rm app}\rangle=|V|.
% $
We define the operator
\[
\mathcal{Q}_\varepsilon(f)
=
\int_V T_\varepsilon(t,x,v')f(v')\,dv'
-
|V|T_\varepsilon(t,x,v)f(v),
\]
and using \eqref{f_expansion} we obtain
\[
(I-\Pi)\mathcal Q_\varepsilon(f^\varepsilon)
=
-\varepsilon\bigl(|V|g^\varepsilon+\rho^\varepsilon\psi_1\bigr)
-\varepsilon^2|V|(I-\Pi)(\psi_1g^\varepsilon).
\]
Applying \(I-\Pi\) {to the kinetic equation \eqref{eq: inital model with scaling}} with
\(\mathcal Q_\varepsilon\) gives
\[
\varepsilon^2 \partial_t g^\varepsilon
+
\varepsilon (I-\Pi)(v\partial_x g^\varepsilon)
+
vM\,\partial_x \rho^\varepsilon
=
-\mu\bigl(|V|g^\varepsilon+\rho^\varepsilon\psi_1\bigr)
-\mu\varepsilon |V|(I-\Pi)(\psi_1g^\varepsilon).
\]
Using \(M=1/|V|\) together with the normalization \(\mu |V|=1\), we obtain
the micro-macro system
\begin{equation}\label{eq:micro-macro-normalized}
\left\{
\begin{aligned}
&\partial_t \rho^\varepsilon
+\partial_x \langle v g^\varepsilon\rangle = 0,\\
&\partial_t g^\varepsilon
+\frac1\varepsilon (I-\Pi)(v\partial_x g^\varepsilon)
=
\frac1{\varepsilon^2}
\left(
-\frac{v}{|V|}\partial_x \rho^\varepsilon
-g^\varepsilon
-\frac{\rho^\varepsilon}{|V|}\psi_1
\right)
-\frac1\varepsilon (I-\Pi)(\psi_1g^\varepsilon).
\end{aligned}
\right.
\end{equation}

}





{
\section{An asymptotic preserving finite difference scheme}
\label{sec:scheme}

In this section we construct an asymptotic preserving finite difference
scheme for the micro-macro formulation \eqref{eq:micro-macro-normalized}. We present the method in one space dimension. The scheme is
designed to remain stable as \(\varepsilon\to0\) and to reduce to a
consistent discretization of the macroscopic drift-diffusion limit.

Let $\Omega = [x_{\min}, x_{\max}]$ be the computational domain and \(x_{j+\frac12}=x_{\min}+(j+\frac12)\Delta x\),
\(x_j=x_{\min}+j\Delta x\), and \(t_n=n\Delta t\). We approximate
\(\rho\), \(S\), and \(N\) at the half grid points, while the microscopic
unknown \(g\) is approximated at the integer grid points,
\[
\rho_{j+\frac12}^n\approx \rho(t_n,x_{j+\frac12}),\qquad
S_{j+\frac12}^n\approx S(t_n,x_{j+\frac12}),\qquad
N_{j+\frac12}^n\approx N(t_n,x_{j+\frac12}),
\]
and $
g_{j,k}^n\approx g(t_n,x_j,v_k),\ v_k\in V_h $ where \(V_h=\{v_k\}\) is a symmetric discrete velocity set satisfying 
$v_{-k}=-v_k$.

For a half-grid scalar variable \(q_{j+\frac12}^n\), we use
\[
\delta_t^+ q_{j+\frac12}^n
=
\frac{q_{j+\frac12}^{n+1}-q_{j+\frac12}^n}{\Delta t},
\ \ \
(\delta_x q)_j^n
=
\frac{q_{j+\frac12}^n-q_{j-\frac12}^n}{\Delta x},\ \ \ \textnormal{and}\ \ \
(\delta_x^2 q)_{j+\frac12}^n
=
\frac{
q_{j+\frac32}^n
-2q_{j+\frac12}^n
+q_{j-\frac12}^n
}{\Delta x^2}.
\]
Moreover, we define the discrete divergence for \(g_{j,k}^n\) by
\[
(\delta_x g)_{j+\frac12,k}^n
=
\frac{g_{j+1,k}^n-g_{j,k}^n}{\Delta x}.
\]
Whenever a half-grid variable is needed at an integer grid point, we use
the local average
\[
q_j^n=\frac12\left(q_{j-\frac12}^n+q_{j+\frac12}^n\right).
\]
The value \((\delta_x^2 q)_j^n\) is defined in the same way by averaging
the two neighboring half-grid values.


We define
\[
\langle \eta\rangle_h=\sum_{k}\omega_k\eta_k,
\qquad \omega_k>0,
\qquad |V|_h=\sum_{k}\omega_k,
\qquad \Pi_h\eta=\frac{\langle\eta\rangle_h}{|V|_h}.
\]
We assume $\langle v\rangle_h=0$. Then $I-\Pi_h$ is the
projection onto the discrete zero-average space.


\paragraph{Discrete response functions}

At each time step, the response functions are evaluated explicitly. 
For $q\in\{S,N\}$, we use the scaled derivative
\begin{equation}
\mathcal D_\varepsilon^v q
=\varepsilon\partial_tq+v\partial_xq.
\end{equation}
We approximate $\partial_tq$ using the equations for \(S\) and
\(N\) at time $t^n$. 
We define
\[
\mathcal L_{S,j}^n
=
\widetilde D_S(\delta_x^2 S)_j^n
-\widetilde\alpha S_j^n
+\widetilde\beta \rho_j^n,
\quad 
\mathcal L_{N,j}^n
=
\widetilde D_N(\delta_x^2 N)_j^n
-\widetilde\gamma \rho_j^n N_j^n .
\]
The discrete response arguments are
\[
\mathcal D_{\varepsilon,h}^{v_k}S_j^n
=\varepsilon\mathcal L_{S,j}^n+v_k(\delta_xS)_j^n,
\qquad
\mathcal D_{\varepsilon,h}^{v_k}N_j^n
=\varepsilon\mathcal L_{N,j}^n+v_k(\delta_xN)_j^n.
\]
All quantities on the right-hand sides are known at time $t^n$.

The velocity average of \(\phi(a+bv)\) does not vanish in general, even
when \(\phi\) is odd. We therefore use the normalized discrete responses
\[
\psi_{S,j,k}^n
=\chi_S\left[
\phi\!\left(\mathcal D_{\varepsilon,h}^{v_k}S_j^n\right)
-\Pi_h\phi\!\left(\mathcal D_{\varepsilon,h}^{v_k}S_j^n\right)
\right],
\psi_{N,j,k}^n
=\chi_N\left[
\phi\!\left(\mathcal D_{\varepsilon,h}^{v_k}N_j^n\right)
-\Pi_h\phi\!\left(\mathcal D_{\varepsilon,h}^{v_k}N_j^n\right)
\right].
\]
We set
$\psi_{j,k}^n
=
\psi_{S,j,k}^n+\psi_{N,j,k}^n ,$
and by construction,
\[
\langle \psi_{S,j,k}^n\rangle_h=0,
\qquad
\langle \psi_{N,j,k}^n\rangle_h=0,
\qquad
\langle \psi_{j,k}^n\rangle_h=0.
\]
Consequently, the discrete turning kernel $
T_{\varepsilon,j,k}^n
=
1+\varepsilon \psi_{j,k}^n$
satisfies
$
\langle T_{\varepsilon,j,k}^n\rangle_h=|V|_h.$
We also introduce
$
B_{S,j}^n
=
\langle v_k\psi_{S,j,k}^n\rangle_h,
\
B_{N,j}^n
=
\langle v_k\psi_{N,j,k}^n\rangle_h.
$
Then, the corresponding discrete drift velocities are defined by
\[
u_{S,j}^n
=
-\frac{B_{S,j}^n}{|V|_h},
\qquad
u_{N,j}^n
=
-\frac{B_{N,j}^n}{|V|_h}.
\]

\paragraph{Fully discrete micro-macro scheme}
The fully discrete AP scheme is given by
\begin{subequations}\label{eq:modified-main-scheme}
\begin{align}
&
\delta_t^+\rho_{j+\frac12}^n
+
\left\langle
v_k(\delta_x g)_{j+\frac12,k}^{n+1}
\right\rangle_h
=0\ ,\label{eq:modified-main-scheme1}
\\[1mm]
&
\delta_t^+ g_{j,k}^n
+
\frac1\varepsilon
(I-\Pi_h)\mathcal A_{j,k}^n
=
\frac1{\varepsilon^2}
\left(
K_{j,k}^{n,n+1}
-
g_{j,k}^{n+1}
\right)
-
\frac1\varepsilon
(I-\Pi_h)(\psi_{j,k}^n g_{j,k}^n)\ ,\label{eq:modified-main-scheme2}
\\[1mm]
&
\delta_t^+ S_{j+1/2}^n
=
\widetilde D_S(\delta_x^2S^{n+1})_{j+1/2}
-\widetilde\alpha S_{j+1/2}^{n+1}
+\widetilde\beta\rho_{j+1/2}^{n+1}\ ,\label{eq:modified-main-scheme3}
\\[1mm]
&
\delta_t^+N_{j+\frac12}^{n}
=
\widetilde D_N(\delta_x^2N^{n+1})_{j+1/2}
-\widetilde\gamma\rho_{j+1/2}^{n+1}N_{j+1/2}^{n+1}\ .\label{eq:modified-main-scheme4}
\end{align}
\end{subequations}
Here the microscopic transport is treated explicitly by the upwind
approximation
\[
\mathcal A_{j,k}^n
=
v_k^+(\delta_x g)_{j-\frac12,k}^{n}
-
v_k^-(\delta_x g)_{j+\frac12,k}^{n},
\]
where
$
v_k^+=\max\{v_k,0\},
\
v_k^-=\max\{-v_k,0\}.
$
The local equilibrium part is defined by
\begin{equation}\label{eq:modified-K-two-upwind}
K_{j,k}^{n,n+1}
=
-\frac{v_k}{|V|_h}(\delta_x\rho)_j^{n+1}
-
\frac{\rho_j^{[S],n+1}}{|V|_h}\psi_{S,j,k}^n
-
\frac{\rho_j^{[N],n+1}}{|V|_h}\psi_{N,j,k}^n .
\end{equation}
The two upwind values $\rho_j^{[S],n+1}$ and $\rho_j^{[N],n+1}$ are chosen according to the effective velocities
\(u_{S,j}^n\) and \(u_{N,j}^n\),
\begin{equation}\label{eq: upwind densities}
\rho_j^{[S],n+1}
=
\begin{cases}
\rho_{j-\frac12}^{n+1}, & u_{S,j}^n\ge 0,\\
\rho_{j+\frac12}^{n+1}, & u_{S,j}^n<0,
\end{cases}
\quad 
\rho_j^{[N],n+1}
=
\begin{cases}
\rho_{j-\frac12}^{n+1}, & u_{N,j}^n\ge 0,\\
\rho_{j+\frac12}^{n+1}, & u_{N,j}^n<0.
\end{cases}
\end{equation}
This choice is consistent with the drift part of the limiting flux since
\[
-\frac{B_{S,j}^n}{|V|_h}\rho_j^{[S],n+1}
-
\frac{B_{N,j}^n}{|V|_h}\rho_j^{[N],n+1}
=
u_{S,j}^n\rho_j^{[S],n+1}
+
u_{N,j}^n\rho_j^{[N],n+1}.
\]
Taking the
discrete velocity average of the second equation in
\eqref{eq:modified-main-scheme} and using
$
\langle K_{j,k}^{n,n+1}\rangle_h=0
$
gives
\[
\langle g_{j,k}^{n+1}\rangle_h
=
\frac{\varepsilon^2}{\varepsilon^2+\Delta t}
\langle g_{j,k}^{n}\rangle_h .
\]
Thus \(\langle g_{j,k}^{n}\rangle_h=0\) is propagated in time if it holds
initially.

\paragraph{Decoupled solution strategy}
Following the micro-macro strategy in \cite{lemou2008new}, the scheme can
be solved in a decoupled way. We first introduce an auxiliary unknown
\(\widetilde g_{j,k}^{n+1}\) by
\begin{equation}\label{eq:modified-gtilda}
\frac{\widetilde g_{j,k}^{n+1}-g_{j,k}^n}{\Delta t}
+
\frac1\varepsilon
(I-\Pi_h)\mathcal A_{j,k}^n
=
\frac1{\varepsilon^2}
\left(
\widetilde K_{j,k}^{n}
-
\widetilde g_{j,k}^{n+1}
\right)
-
\frac1\varepsilon
(I-\Pi_h)(\psi_{j,k}^n g_{j,k}^n).
\end{equation}
where
\[
\widetilde K_{j,k}^{n}
=
-\frac{v_k}{|V|_h}(\delta_x\rho)_j^{n}
-
\frac{\rho_j^{[S],n}}{|V|_h}\psi_{S,j,k}^n
-
\frac{\rho_j^{[N],n}}{|V|_h}\psi_{N,j,k}^n .
\]
The upwind values \(\rho_j^{[S],n}\) and \(\rho_j^{[N],n}\) are defined in
the same way as in \eqref{eq: upwind densities}, with \(\rho^{n+1}\) replaced by \(\rho^n\).

All terms in \eqref{eq:modified-gtilda} are known except
\(\widetilde g_{j,k}^{n+1}\). Hence
\[
\widetilde g_{j,k}^{n+1}
=
\frac{\varepsilon^2}{\varepsilon^2+\Delta t}g_{j,k}^n
-
\frac{\varepsilon\Delta t}{\varepsilon^2+\Delta t}
(I-\Pi_h)\mathcal A_{j,k}^n
+
\frac{\Delta t}{\varepsilon^2+\Delta t}
\widetilde K_{j,k}^{n}
-
\frac{\varepsilon\Delta t}{\varepsilon^2+\Delta t}
(I-\Pi_h)(\psi_{j,k}^n g_{j,k}^n).
\]
Comparing \eqref{eq:modified-gtilda} with 
\eqref{eq:modified-main-scheme2}, we obtain
\begin{equation}\label{eq:modified-g-reconstruction}
g_{j,k}^{n+1}
=
\widetilde g_{j,k}^{n+1}
+
\lambda_\varepsilon
\left(
K_{j,k}^{n,n+1}
-
\widetilde K_{j,k}^{n}
\right),
\end{equation}
where
\begin{equation}\label{eq: lambda_eps}
\lambda_\varepsilon
=
\frac{\Delta t}{\varepsilon^2+\Delta t}.
\end{equation}

Substituting \eqref{eq:modified-g-reconstruction} into  \eqref{eq:modified-main-scheme1}, we obtain a closed equation for
\(\rho^{n+1}\). Define
\[
D_{\rho,h}
=
\frac{\langle v_k^2\rangle_h}{|V|_h},
\qquad
b_{S,j}^n
=
\frac{B_{S,j}^n}{|V|_h},
\qquad
b_{N,j}^n
=
\frac{B_{N,j}^n}{|V|_h},
\]
and let
$
a^n=\lambda_\varepsilon D_{\rho,h},
\
c_{S,j}^n=\lambda_\varepsilon b_{S,j}^n,
\
c_{N,j}^n=\lambda_\varepsilon b_{N,j}^n.
$
Then the density equation can be written as
\begin{equation}\label{eq:modified-rho-equation-two-upwind}
\frac{\rho_{j+\frac12}^{n+1}}{\Delta t}
-
\delta_x\left(
a^n\delta_x\rho^{n+1}
\right)_{j+\frac12}
-
\delta_x\left(
c_S^n\rho^{[S],n+1}
\right)_{j+\frac12}
-
\delta_x\left(
c_N^n\rho^{[N],n+1}
\right)_{j+\frac12}
=
r_{j+\frac12}^n,
\end{equation}
where the known right-hand side is
\[
\begin{aligned}
r_{j+\frac12}^n
=&
\frac{\rho_{j+\frac12}^n}{\Delta t}
-
\left\langle
v_k(\delta_x\widetilde g)_{j+\frac12,k}^{n+1}
\right\rangle_h
-
\delta_x\left(
a^n\delta_x\rho^n
\right)_{j+\frac12}
\\ &{}-
\delta_x\left(
c_S^n\rho^{[S],n}
\right)_{j+\frac12}
-
\delta_x\left(
c_N^n\rho^{[N],n}
\right)_{j+\frac12}.
\end{aligned}
\]
Thus only a scalar linear system for \(\rho^{n+1}\) needs to be solved at
each time step. With the upwind choice for the drift terms \eqref{eq: upwind densities}, the
resulting matrix is tridiagonal. Once \(\rho^{n+1}\) is obtained,
\(g^{n+1}\) is recovered explicitly from
\eqref{eq:modified-g-reconstruction}. 
%Finally, \(S^{n+1}\) and \(N^{n+1}\) are updated by  \eqref{eq:modified-main-scheme3} and \eqref{eq:modified-main-scheme4}.
We impose zero numerical density flux at both boundaries.
This condition is preserved by the positivity correction
described in Appendix~\ref{sec:limiter}.
The corrected density is then used to update $S^{n+1}$ and $N^{n+1}$
independently by \eqref{eq:modified-main-scheme3} and
\eqref{eq:modified-main-scheme4}, respectively.



\subsection{Asymptotic preserving property}

We now verify that the scheme is formally asymptotic preserving
Let \(\Delta t\) and \(\Delta x\) be fixed. Notice that $\lambda_\varepsilon$ in \eqref{eq: lambda_eps}
 converges to \(1\) as \(\varepsilon\) goes to zero.
Moreover, by the definition of \(\widetilde g^{n+1}\),
\(\widetilde g_{j,k}^{n+1}\) converges to \(\widetilde K_{j,k}^{n}\).
Using the expression of \(\widetilde K^n\), we obtain \(\left\langle
v_k(\delta_x\widetilde g)_{j+\frac12,k}^{n+1}
\right\rangle_h\) converges to
\[
-\delta_x\left(D_{\rho,h}\delta_x\rho^n\right)_{j+\frac12}
-
\delta_x\left(b_S^n\rho^{[S],n}\right)_{j+\frac12}
-
\delta_x\left(b_N^n\rho^{[N],n}\right)_{j+\frac12}.
\]
It follows from the definition of \(r_{j+\frac12}^n\) in
\eqref{eq:modified-rho-equation-two-upwind} that
\(
r_{j+\frac12}^n
\quad\text{converges to}\quad
\frac{\rho_{j+\frac12}^n}{\Delta t}.
\)
Therefore, in the limit as \(\varepsilon\) goes to zero,
\eqref{eq:modified-rho-equation-two-upwind} reduces to
\begin{equation}\label{eq:limiting-density-scheme}
\frac{\rho_{j+\frac12}^{n+1}-\rho_{j+\frac12}^{n}}{\Delta t}
-
\delta_x\left(
D_{\rho,h}\delta_x\rho^{n+1}
\right)_{j+\frac12}
-
\delta_x\left(
b_S^n\rho^{[S],n+1}
\right)_{j+\frac12}
-
\delta_x\left(
b_N^n\rho^{[N],n+1}
\right)_{j+\frac12}
=0 .
\end{equation}
This is a consistent implicit finite difference discretization of the
limiting drift-diffusion equation
\[
\partial_t\rho
-
\partial_x\left(D_{\rho,h}\partial_x\rho\right)
-
\partial_x\left(b_S\rho\right)
-
\partial_x\left(b_N\rho\right)
=0\ ,
\]
where \(u_S=-b_S\) and \(u_N=-b_N\).
Hence, for fixed $\Delta t$ and $\Delta x$, the fully discrete scheme reduces to a consistent discretization of the macroscopic drift–diffusion limit as $\varepsilon\to0$, which is the formal AP property in the diffusive regime.

\subsection{Bound-preserving property}

We now discuss the bound-preserving property of the limiting scheme.
In general, the full kinetic scheme is not automatically positivity
preserving. In the computation, this issue is handled by a post-processing
limiter as explained in Appendix \ref{sec:limiter}. For the limiting macroscopic scheme, however, a positivity and
boundedness result can be established.

In this subsection, we consider the limiting macroscopic scheme \eqref{eq:limiting-density-scheme} with
homogeneous no-flux boundary conditions, so that the discrete boundary
fluxes vanish. 

% Recall that the limiting scheme reads
% \[
% \delta_t^+\rho_{j+\frac12}^n
% =
% D_{\rho,h}\delta_x^2\rho_{j+\frac12}^{n+1}
% -
% \delta_x\left(u_S^n\rho^{[S],n+1}\right)_{j+\frac12}
% -
% \delta_x\left(u_N^n\rho^{[N],n+1}\right)_{j+\frac12},
% \]
% where the upwind values are chosen according to $u_{S,j}^n$ and $u_{N,j}^n$ as in \eqref{eq: upwind densities}

The limiting scheme can be written in matrix form as
\[
M^n\brho^{n+1}=\brho^n \ ,
\]
where
 $\boldsymbol{\rho}^{\,n}
=
(\rho_{j+\frac12}^n)_j$
denotes the vector of half-grid density values at time level $n$.
The tridiagonal matrix $M^n$ has entries
\begin{align*}
&m_{j,j-1}^n
=
-D_{\rho,h}\frac{\Delta t}{(\Delta x)^2}
-\frac{\Delta t}{\Delta x}(u_{S,j}^n)^+
-\frac{\Delta t}{\Delta x}(u_{N,j}^n)^+,\\
&m_{j,j}^n
=
1+2D_{\rho,h}\frac{\Delta t}{(\Delta x)^2}
+\frac{\Delta t}{\Delta x}
\left[(u_{S,j+1}^n)^+ +(u_{S,j}^n)^-\right]
+\frac{\Delta t}{\Delta x}
\left[(u_{N,j+1}^n)^+ +(u_{N,j}^n)^-\right],\\
&m_{j,j+1}^n
=
-D_{\rho,h}\frac{\Delta t}{(\Delta x)^2}
-\frac{\Delta t}{\Delta x}(u_{S,j+1}^n)^-
-\frac{\Delta t}{\Delta x}(u_{N,j+1}^n)^- .
\end{align*}
Here $(a)^+=\max\{a,0\}$ and $(a)^-=\max\{-a,0\}$.
The boundary rows of \(M^n\) are modified according to the homogeneous
no-flux boundary conditions.


\begin{proposition}[Mass conservation, positivity, and upper bound]
Assume that $D_{\rho,h}>0$ and that the effective velocities satisfy
\[
|u_{S,j}^n|\le U_S,\qquad |u_{N,j}^n|\le U_N
\]
for all $j$ and $n$. Suppose further that
\begin{equation}\label{eq:CFL-bound-new}
\frac{\Delta t}{\Delta x}
\le
\frac{1}{4(U_S+U_N)} .
\end{equation}
If the initial data satisfy $\rho_{j+\frac12}^0\ge0$ for all $j$, then the
limiting scheme satisfies
\[
\rho_{j+\frac12}^n\ge0,\qquad \forall j,\ n\ge0.
\]
Moreover, under the homogeneous no-flux boundary conditions, the discrete
mass is conserved,
\[
\|\brho^{n+1}\|_1=\|\brho^n\|_1 := \sum_j|\rho^n_{j+\frac12}|,
\]
and the solution satisfies the uniform bound in time on a fixed spatial grid
\[
0\le \rho_{j+\frac12}^n\le \bar\rho,\qquad
\bar\rho=2\|\brho^0\|_1 .
\]
\end{proposition}

\begin{proof}{\rm
We first observe that
\[
m_{j,j-1}^n+m_{j,j}^n+m_{j,j+1}^n
=
1+\frac{\Delta t}{\Delta x}
\left(u_{S,j+1}^n-u_{S,j}^n\right)
+\frac{\Delta t}{\Delta x}
\left(u_{N,j+1}^n-u_{N,j}^n\right)
=:F_j^n .
\]
Using the bounds on $u_S^n$ and $u_N^n$, we have
\[
F_j^n
\ge
1-2(U_S+U_N)\frac{\Delta t}{\Delta x}.
\]
Under \eqref{eq:CFL-bound-new}, this gives
$
F_j^n\ge \frac12 .
$

The off-diagonal entries of $M^n$ are nonpositive. Moreover, the above
lower bound implies the required diagonal dominance under the CFL
condition. Hence $M^n$ is a nonsingular M-matrix and $(M^n)^{-1}$ is
nonnegative. Therefore, if $\brho^n\ge0$, then
\[
\brho^{n+1}=(M^n)^{-1}\brho^n\ge0.
\]
Mass conservation follows from the conservative form of the scheme and the
homogeneous no-flux boundary conditions.  More precisely, 
\[
\|\brho^{n+1}\|_1
=
\mathbf 1^T\brho^{n+1}
=
\mathbf 1^T M^n\brho^{n+1}
=
\mathbf 1^T\brho^n
=
\|\brho^n\|_1 ,
\]
where we used the relation
\(
\mathbf 1^T M^n=\mathbf 1^T .
\)

Finally, let $\bar\rho=2\|\brho^0\|_1$. By mass conservation and
positivity,
\[
\rho_{j+\frac12}^n\le \|\brho^n\|_1=\|\brho^0\|_1=\frac12\bar\rho .
\]
Since $F_j^n\ge1/2$, we obtain
\[
(M^n(\bar\rho\mathbf 1))_j
=
F_j^n\bar\rho
\ge
\frac12\bar\rho
\ge
\rho_{j+\frac12}^n .
\]
Using the nonnegativity of $(M^n)^{-1}$, we conclude that
\[
\brho^{n+1}
=
(M^n)^{-1}\brho^n
\le
(M^n)^{-1}M^n(\bar\rho\mathbf 1)
=
\bar\rho\mathbf 1 .
\]
This proves the desired upper bound. }
\end{proof}



\section{Numerical experiments}\label{sec: numerics}

In this section, we present numerical experiments to illustrate the
traveling pulse behavior of the model and to verify the asymptotic
preserving property of the proposed micro-macro scheme. We first validate
the macroscopic solver by comparing the numerical traveling pulse with
the analytical traveling profile. We then study the dependence of the
traveling speed on the chemotactic sensitivities. Finally, we compare
the kinetic and macroscopic solutions for different values of
\(\varepsilon\).


Unless otherwise stated, we consider the one-dimensional computational
domain \(\Omega=(0,150)\). For the profile and asymptotic preserving
tests, we use
\begingroup\color{red}
\[
\Delta x=0.025,
\qquad
\Delta t=10^{-3}.
\]
\endgroup
The velocity domain is \(V=[-1,1]\). It is discretized by a uniform
trapezoidal quadrature with \(\Delta v=0.02\). We denote the discrete
velocity nodes and weights by \(\{v_k,\omega_k\}\), and use
\[
\langle h\rangle_h=\sum_k\omega_k h(v_k),
\qquad
|V|_h=\sum_k\omega_k .
\]
The discrete diffusion coefficient in the limiting macroscopic equation
is
 $D_{\rho,h}
=
\frac{\langle v^2\rangle_h}{|V|_h}.$
For the default velocity grid with \(\Delta v=0.02\), this gives
\(D_{\rho,h}=0.3334\).


We use the reflecting boundary conditions described in
Appendix~\ref{sec:limiter} for the kinetic variable and
homogeneous Neumann conditions for the chemical variables.


The initial data are chosen so that cells are concentrated near the left
side of the domain, with no chemoattractant signal and a constant
nutrient level. More precisely, we take
\[
\rho(0,x)=\rho_0 e^{-\lambda x},
\qquad
S(0,x)=0,
\qquad
N(0,x)=N_0,
\]
where
$\rho_0=1,
\
\lambda=3,
\
N_0=10^3 .$
For the kinetic computations, the initial data are chosen consistently
with the micro-macro decomposition,
\[
f^\varepsilon(0,x,v)
=
\frac{\rho(0,x)}{|V|_h}
+
\varepsilon g(0,x,v),
\qquad
g(0,x,v)=0 .
\]
The chemical parameters are
$
D_S=2,
\
D_N=0,
\
\alpha=0.05,
\
\beta=\gamma=1 .
$
Unless otherwise stated, the chemotactic sensitivities in the profile
and AP tests are
$\chi_S=0.5,
\
\chi_N=1.1 .$
% REVIEW: sign-response-normalization
\begin{resultrevision}
For the sign-response experiments in Sections~\ref{sec:macro-validation}--\ref{sec:sign-ap-experiments}, we use
\[
\phi_h(y)=-s_h\operatorname{sign}(y),\qquad
s_h=\frac{|V|_h}{\langle|v|\rangle_h}.
\]
This normalization gives the macroscopic drift velocities
$u_S=\chi_S\operatorname{sign}(S_x)$ and
$u_N=\chi_N\operatorname{sign}(N_x)$.
The response arguments use $\varepsilon=0$ for the macroscopic model
and the positive kinetic $\varepsilon$ for the AP model.
The kinetic computations use the correction in Appendix~\ref{sec:limiter}.
\end{resultrevision}

\subsection{Macroscopic traveling pulse and analytical profile}\label{sec:macro-validation}

We first validate the macroscopic solver by comparing the numerical
traveling pulse with the analytical traveling profile. Let
\(x_p\) be the peak position of the numerical density and set
\[
z=x-x_p,
\qquad
z_-=\min\{z,0\},
\qquad
z_+=\max\{z,0\}.
\]
The analytical traveling profile with the response function $\phi(y)=-\operatorname{sign}(y)$ found in \cite{calvez2019chemotactic} is written as 
\begin{equation}
\rho_{\rm tw}(x)
=
\rho_p
\exp\left(
\lambda_L z_-
+
\lambda_R z_+
\right),
\label{eq:analytic-sol}
\end{equation}
where
\[
\lambda_L
=
\frac{\chi_S+\chi_N-\sigma}{D_{\rho,h}},
\qquad
\lambda_R
=
\frac{\chi_N-\chi_S-\sigma}{D_{\rho,h}} .
\]
Here \(\lambda_L>0\) and \(\lambda_R<0\), so that the profile decays on
both sides of the peak. The amplitude \(\rho_p\) is chosen by discrete
mass matching. More precisely, if \(z_j=x_{j+\frac12}-x_p\), then
\[
\rho_p
=
\frac{
\Delta x\sum_j \rho_j
}{
\Delta x\sum_j
\exp\left(
\lambda_L (z_j)_-
+
\lambda_R (z_j)_+
\right)
}.
\]

% REVIEW: figure1-profile-comparison
\begin{resultrevision}
Figure~\ref{fig:macro_analytic_cmp} compares the numerical macroscopic pulse
with the analytical profile at $t=30$ and $t=50$.
For this macroscopic profile comparison, we use $\Delta x=0.025$ and
$\Delta t=0.001$.
The numerical profiles reproduce the asymmetric analytical shape,
with peak-height discrepancies of approximately $4$--$6\%$.
The peaks are aligned as defined above, so this compares profile shape
rather than independently testing propagation speed.
The AP experiments below now use this same spatial grid and physical
parameters; their macroscopic reference is recomputed independently.
\end{resultrevision}

\begin{figure}[htbp]
    \centering
    \begin{subfigure}{0.48\linewidth}
        \centering
        \includegraphics[width=\linewidth]{macro_compare_t_030p000.eps}
        \caption{\(t=30\)}
        \label{fig:macro_analytic_t30}
    \end{subfigure}
    \begin{subfigure}{0.48\linewidth}
        \centering
        \includegraphics[width=\linewidth]{macro_compare_t_050p000.eps}
        \caption{\(t=50\)}
        \label{fig:macro_analytic_t50}
    \end{subfigure}
    \caption{Comparison between the macroscopic numerical profiles and
    the analytical traveling profiles at \(t=30\) and \(t=50\). }
    \label{fig:macro_analytic_cmp}
\end{figure}

\subsection{Dependence of the traveling speed on chemotactic sensitivities}\label{sec:sensitivity-study}

We next investigate how the traveling speed depends on \(\chi_S\) and
\(\chi_N\). For each pair \((\chi_S,\chi_N)\), the predicted traveling
speed \(\sigma\) is determined by
\begin{equation}
\chi_N-\sigma
=
\chi_S
\frac{\sigma}{\sqrt{4D_S\alpha+\sigma^2}} .
\label{eq:sigma-macro}
\end{equation}
To estimate the traveling speed numerically, we track the center of
mass of the density,
\begin{equation}\label{eq: center of mass}
X_\rho(t_n)
=
\frac{
\Delta x\sum_j x_{j+\frac12}\rho_{j+\frac12}^n
}{
\Delta x\sum_j \rho_{j+\frac12}^n
}.
\end{equation}


% REVIEW: figure2-source-and-interpretation
\begin{resultrevision}
Figure~\ref{fig:traveling_speed} retains the original sensitivity-scan data,
with $\chi_S$ and $\chi_N$ sampled on $[0,2]\times[0,2]$.
The traced macroscopic data use $\Omega=(0,300)$, $\Delta x=0.1$,
$\Delta t=10^{-2}$ and $T=100$.
The kinetic panel at $\varepsilon=0.01$ is retained as a qualitative
comparison because its complete discretization settings could not be verified.
The panels illustrate the broad dependence on the competing sensitivities;
they are not used as a quantitative validation of the corrected solver.
The updated convergence tests below instead use the common baseline
settings stated at the beginning of this section.
\end{resultrevision}

\begin{figure}[htbp]
    \centering
    \begin{subfigure}{0.3\linewidth}
        \centering
        \includegraphics[width=\linewidth]{travelling_speed_analytic.eps}
        \caption{}
        \label{fig:analytic_sigma}
    \end{subfigure}
    \begin{subfigure}{0.3\linewidth}
        \centering
        \includegraphics[width=\linewidth]{travelling_speed_numeric_macro.eps}
        \caption{}
        \label{fig:numerical_macro_sigma}
    \end{subfigure}
    \begin{subfigure}{0.3\linewidth}
        \centering
        \includegraphics[width=\linewidth]{travelling_speed_numeric_kinetic.eps}
        \caption{}
        \label{fig:numerical_kinetic_sigma}
    \end{subfigure}
    \caption{\st{Dependence of the traveling speed on $\chi_S$ and $\chi_N$.
    Panel (a) shows the speed predicted by \eqref{eq:sigma-macro};
    panels (b) and (c) retain the earlier macroscopic and kinetic sensitivity scans,
    with $\varepsilon=0.01$ in panel (c), for qualitative comparison.}}
    \label{fig:traveling_speed}
\end{figure}

\subsection{Asymptotic preserving convergence}\label{sec:sign-ap-experiments}

We now compare the kinetic solution with the limiting macroscopic
solution for different values of \(\varepsilon\). The purpose of this
test is to verify that, for fixed \(\Delta x\), \(\Delta v\), and
\(\Delta t\), the kinetic AP solution approaches the limiting
macroscopic solution as \(\varepsilon\) decreases. We use the baseline
parameter setting described at the beginning of this section and consider
$\varepsilon=0.1,
\
\varepsilon=0.05,
\
\varepsilon=0.01.$

% REVIEW: figure3-density-comparison
\begin{resultrevision}
Figure~\ref{fig:conv_T} compares the kinetic and numerical macroscopic
density pulses at six times. The profiles become closer as
$\varepsilon$ decreases, although finite differences remain.
The displayed interval $0\le x\le40$ is a profile window within the
full computational domain; at late times the pulse tails extend beyond
its right endpoint. All density errors below are evaluated on the
complete domain without shifting or rescaling the profiles.
\end{resultrevision}

\begin{figure}[htbp]
    \centering
    \includegraphics[width=\linewidth]{AP_profiles.eps}
    \caption{\st{Kinetic and numerical macroscopic density pulses at
    $t=1,10,20,30,40,50$, shown in the window $0\le x\le40$.
    Late-time tails extend beyond this window; the errors in
    Figure~\ref{fig:conv_error} use the complete domain.}}
    \label{fig:conv_T}
\end{figure}

To quantify the convergence, we use the discrete \(L^2\)-error
\[
\operatorname{err}_{\rho,\varepsilon}(t)
=
\left(
\Delta x\sum_j
\left|
\rho_{j+\frac12}^{\varepsilon}(t)
-
\rho_{j+\frac12}(t)
\right|^2
\right)^{1/2},
\]
where \(\rho^\varepsilon\) is reconstructed from the kinetic solution
and \(\rho\) denotes the solution of the limiting macroscopic model.
% REVIEW: figure4-full-domain-errors
\begin{resultrevision}
Figure~\ref{fig:conv_error} shows the full-domain error at $T=50$. Across the seven sampled values, it decreases from $4.7168\times10^{-2}$ at $\varepsilon=0.1$ to $2.2208\times10^{-4}$ at $\varepsilon=0.001$. This supports convergence toward the numerical macroscopic solution on the common fine grid. The reference lines are comparisons, not fitted convergence rates.
\end{resultrevision}

\begin{figure}[htbp]
    \centering
    \includegraphics[width=0.6\linewidth]{AP_error_T50.eps}
    \caption{Final-time density error
    \(\operatorname{err}_{\rho,\varepsilon}(50)\) versus
    \(\varepsilon\). The reference lines \(O(\varepsilon)\) and
    \(O(\varepsilon^2)\) are included for comparison.}
    \label{fig:conv_error}
\end{figure}

% REVIEW: figure5-finite-time-speeds
\begin{resultrevision}
Figure~\ref{fig:conv_speed} compares speeds estimated from the center-of-mass trajectories in \eqref{eq: center of mass}. The dashed line is the sign-response prediction $\sigma_{\rm pred}\simeq0.724$ from \eqref{eq:sigma-macro}. Linear fits on $[40,50]$ give $0.73101$, $0.72483$ and $0.72331$ for $\varepsilon=0.1,0.05,0.01$, respectively, compared with $0.72316$ for the numerical macroscopic solution. The speed discrepancy decreases over these three sampled values. This finite-time comparison does not establish long-time stability of a traveling wave.
\end{resultrevision}

\begin{figure}[htbp]
    \centering
    \includegraphics[width=0.82\linewidth]{AP_speed_direct.eps}
    \caption{Traveling speed comparison between the macroscopic model
    and the kinetic model for different values of \(\varepsilon\). 
    The horizontal dashed line indicates the macroscopic predicted speed \(\sigma_{\rm macro}\approx0.724\), obtained from \eqref{eq:sigma-macro}.}
    \label{fig:conv_speed}
\end{figure}

\begin{samepage}
% REVIEW: sign-results-summary
\begin{resultrevision}
Together, the full-domain density errors and the propagation-speed
comparisons support the approach of the kinetic solution to the
numerical macroscopic limit for the tested sign-response cases.
These are finite-grid, finite-time observations rather than a uniform
error estimate or a stability result for traveling waves.
\end{resultrevision}
\end{samepage}

\FloatBarrier

}


% This subsection is locally red for review; curve colors are unchanged.
\FloatBarrier
\begin{resultrevision}
\subsection{Traveling pulses with an exponential response}\label{sec:exponential-response}

We next use the exponentially modulated response
\begin{equation}\label{eq:exponential-response}
\phi_{a,h}(y)=-s_h e^{-a|y|}\operatorname{sign}(y),\qquad
s_h=\frac{|V|_h}{\langle|v|\rangle_h}.
\end{equation}
The parameter $a\ge0$ is distinct from the chemical degradation rate $\alpha=0.05$;
$a=0$ recovers the normalized sign response exactly.
The remaining settings follow Section~\ref{sec:sign-ap-experiments}.
Both chemical responses retain the full kinetic argument and velocity centering.
The macroscopic drift uses weighted moments of the same response at $\varepsilon=0$.

For $a=0.001$ and $\varepsilon=0.01$, Figure~\ref{fig:exponential-propagation} shows a propagating pulse with limited late-time deformation. For the saved profiles at $t=30,40,50$, centroid-translated shapes differ by up to $13.4\%$ in relative $L^2$. The centroid slopes on $[30,40]$ and $[40,50]$ differ by $0.08\%$. At fixed $a$, the full-domain density errors relative to the corresponding numerical macroscopic solution decrease with $\varepsilon$ in Figure~\ref{fig:exponential-limits}(a). At fixed $\varepsilon=0.01$, the density discrepancies from the numerical sign reference decrease with $a$, as illustrated in Figure~\ref{fig:exponential-limits}(b). The fitted propagation speeds also approach the computed sign value as $a$ decreases (Table~\ref{tab:exponential-speeds}). These comparisons retain physical positions and amplitudes; they describe finite-time behavior on the tested grid, without establishing long-time stability.

\begin{figure}[htbp]
\centering
\includegraphics[width=\linewidth]{exponential_pulse_evolution.eps}
\captionsetup{labelfont={color=red},textfont={color=red}}
\caption{Density evolution with the exponential response at $a=0.001$ and $\varepsilon=0.01$.
Densities at $t=1,10,20,30,40,50$ retain their original physical positions and amplitudes.
The plotting window is $0\le x\le40$; late-time tails extend beyond its right endpoint.}
\label{fig:exponential-propagation}
\end{figure}

\begin{figure}[htbp]
\centering
\includegraphics[width=\linewidth]{exponential_two_limits.eps}
\captionsetup{labelfont={color=red},textfont={color=red}}
\caption{Density comparisons at $T=50$.
(a) Fixed $a=0.001$: kinetic densities for $\varepsilon=0.1,0.05,0.01$
and the corresponding numerical macroscopic density (black solid line).
(b) Fixed $\varepsilon=0.01$: densities for $a=0.001,0.0005,0.0001,0$;
the sign reference $a=0$ is black and solid.
No translation or amplitude normalization is applied.}
\label{fig:exponential-limits}
\end{figure}

\begin{table}[htbp]
\centering
\color{red}
\captionsetup{labelfont={color=red},textfont={color=red}}
\caption{Finite-time numerical propagation speeds for the exponential response.}
\label{tab:exponential-speeds}
\begin{tabular}{lcccc}
\hline
$a$ & 0.001 & 0.0005 & 0.0001 & 0 \\
\hline
Numerical propagation speed & 0.690 & 0.707 & 0.720 & 0.724 \\
\hline
\end{tabular}
\par\smallskip
\begin{minipage}{0.92\linewidth}
\color{red}\footnotesize
Values are slopes of least-squares fits to the center of mass in \eqref{eq: center of mass}
over $t\in[30,50]$, with $\varepsilon=0.01$.
They estimate propagation on this interval, without asserting an exact steady traveling profile.
\end{minipage}
\end{table}
\FloatBarrier
\end{resultrevision}


\section{Conclusions}\label{sec: conclusions}

In this work, we developed and analyzed an asymptotic preserving finite difference scheme for a kinetic chemotaxis model describing bacterial traveling pulses. The model couples a run-and-tumble kinetic equation with reaction--diffusion equations for a nutrient and a self-produced chemoattractant, motivated by experimental observations of \textit{E.~coli} pulse propagation in microchannels.

The proposed numerical method is based on a micro--macro decomposition of the kinetic equation and remains stable in both the kinetic and diffusive regimes. In the limit as the scaling parameter $\varepsilon \to 0$, the scheme naturally reduces to a consistent discretization of the limiting drift--diffusion system. We also established positivity, mass conservation, and uniform boundedness properties for the limiting macroscopic scheme, and introduced a positivity-preserving post-processing procedure for the kinetic solver.

% REVIEW: conclusion-numerical-results
\begin{resultrevision}
The sign-response experiments produce propagating asymmetric pulses.
The numerical profiles retain visible discretization differences from
the mass-matched analytical profile, while their late-time speeds are
close to the predicted value. On the tested fixed grid, the full-domain
density errors decrease as $\varepsilon$ decreases, and the three
reported kinetic speeds approach the numerical macroscopic speed.
These comparisons support the observed diffusive-limit behavior
without asserting an exact profile match or long-time wave stability.
\end{resultrevision}

Several directions remain open for future work. A rigorous convergence analysis of the fully discrete AP scheme would be of significant interest. It would also be natural to extend the method to higher-dimensional settings and more general kinetic response functions. Another important direction concerns the development of fully positivity-preserving AP schemes without post-processing corrections. Finally, incorporating additional biological mechanisms, such as population growth, signal adaptation, or interactions between multiple bacterial species, could provide a more realistic description of collective bacterial migration phenomena.



\bigskip
\noindent{\large\bfseries Acknowledgements} 
CA was supported by TÜBİTAK (The Scientific and Technological Research Council of Türkiye) through the International Postdoctoral Research Fellowship Program under Grant No: 1059B191800629.
JAC was supported by the Advanced Grant Nonlocal-CPD (Nonlocal PDEs for Complex Particle Dynamics: Phase Transitions, Patterns and Synchronization) of the European Research Council Executive Agency (ERC) under the European Union Horizon 2020 research and innovation programme (grant agreement No. 883363) and  partially supported by the EPSRC EP/V051121/1. JAC was partially supported by the ``Maria de Maeztu'' Excellence Unit IMAG, reference CEX2020-001105-M, funded by MCIN/AEI/10.13039/501100011033/.
GER is supported by the Spanish grant PID2022-143012NA-100 and she is a member of the Catalan research group 2021-SGR-00087.
XR acknowledges support from the R\&D Program of Beijing Municipal Education Commission from China under Grant No.  KM202310028016, and the National Natural Science Foundation of China under Grant No. 12201436. 



\appendix

% \section{Numerical scheme for the PDE}
% For the numerical solution of the PDE we rewrite \eqref{eq: rho nondimension} as
% \begin{equation}
%     \partial_t\rho=\nabla\cdot\Bigl(\rho(A\nabla H'(\rho)-Bu_S-Cu_N)\Bigr)\ .
% \end{equation}
% % \textcolor{red}{(Explain numerical discretisation)}

% Via standard upwind finite difference discretization, we have
% \begin{equation}
% \frac{d\rho_j}{dt}=-\frac{F_{j+\frac{1}{2}}(t)-F_{j-\frac{1}{2}}(t)}{\Delta x},
% \end{equation}
% where
% \begin{equation}
% F_{j+\frac{1}{2}}=U^+_{j+\frac{1}{2}}\rho_j+U^-_{j+\frac{1}{2}}\rho_{j+1}, \quad  F_{j-\frac{1}{2}}=U^+_{j-\frac{1}{2}}\rho_{j-1}+U^-_{j-\frac{1}{2}}\rho_j\ .
% \end{equation}
% Also
% \[
% \xi_i=AH'(\rho)=A\log(\rho_i)
% \]
% and
% \begin{equation}
%     U_{j+\frac{1}{2}}=-\frac{\xi_{j+1}-\xi_j}{\Delta x}+a_j[S]+a_j[N]\ ,\ \quad    U_{j-\frac{1}{2}}=-\frac{\xi_{j}-\xi_{j-1}}{\Delta x}+a_{j-1}[S]+a_{j-1}[N]\ ,
% \end{equation}
% \[
% a_j[S]=B\sign\Bigl( \frac{S_{j+1}-S_j}{\Delta x}\Bigr)\ ,\ \quad a_{j-1}[S]=B\sign\Bigl( \frac{S_{j}-S_{j-1}}{\Delta x}\Bigr)
% \]
% and finally
% \begin{equation}
%     U_{j+\frac{1}{2}}^+=\max(U_{j+\frac{1}{2}},0)\ ,\quad U_{j+\frac{1}{2}}^-=\min(U_{j+\frac{1}{2}},0)
% \end{equation}

% We consider zero-flux boundary conditions for $\rho$, $S$ and $N$.


\iffalse
\section{Tumbling operator}\label{app: T}

\begin{figure}[!h]
    \centering
    \includegraphics[width=0.6\linewidth]{tumbling rates.png}
    \caption{In this figure there are plotted the profiles of $S$ and $N$ on the plane defined by the axes $z=0$ and $v=c$. In each quadrant we have the corresponding value of the tumbling rate.}
    \label{tumbling rates}
\end{figure}

To better understand these values, we will explain the intuition behind the signs in each expression of $T$.

\begin{itemize}
    \item $T^+_- = 1-\chi_S-\chi_N$: In this particular case, the $+$ tells us that $v>c$, i.e. the bacteria of this group are moving faster than the wave, while the $-$ corresponds to $z<0$ which means that bacteria are at the left side of the peak of $S$ concentration. Therefore we are in the situation where bacteria move toward the maximum concentration of $S$ and also toward higher concentration of $N$, since it is a favorable motion we do not want to change it that much so we need the tumbling rate as small as possible, so $T^+_-$ is minimum.
    \item $T^+_+ = 1+\chi_S-\chi_N$: Here, we also have $v>c$ which indicates that the bacteria are moving faster than the wave (to the right). As far as the nutrient is concerned, we are going in the right direction, thus, $\chi_N$ does not contribute to more tumbling, on the contrary, wants less tumbling, and that is why we have $-\chi_N$ as well as in the case above. On the other hand, the subscript $+$ shows that $z>0$, i.e. we are in the right side of the peak of $S$. So, as far as the $S$ is concerned, bacteria move away from the higher concentration of $S$ (not favorable for $S$). Consequently, the factor $\chi_S$ does contribute to a higher tumbling rate to revert this situation, that is why we have a positive $+\chi_S$.

    \item $T^-_- = 1+\chi_S+\chi_N$: In the third setup we have $v<c$ which means that we are moving slower or even in the opposite direction to the wave, and $z<0$ indicating bacteria are in the left side of the peak of $S$. We are in the worst scenario because bacteria move away (or goes slowly) toward the peak of $S$ and from the higher concentration of $N$ too. Taking this into account the best is to tumble more to see if we can rectify this situation. Hence, the $T$ is maximal: both factors contribute to have more tumbling.

    \item $T^-_+ = 1-\chi_S+\chi_N$: Finally, we have the setup where the $-\chi_S$ does not contribute to more tumbling because, since bacteria are at the right side of the peak, moving backwards or slower it is fine for $S$. Oppositely, it is not a good situation for $N$, so the factor $+\chi_N$ does contribute to more tumbling.
\end{itemize}
\fi

\section{Boundary treatment and positivity preservation}
\label{sec:limiter}
We use the discrete velocity average
\[
\langle h\rangle_h=\sum_{k}\omega_k h(v_k),
\qquad \omega_k>0,\qquad
|V|_h=\sum_{k}\omega_k,\qquad
\psi_0(v_k)=|V|_h^{-1}.
\]
Let $J$ be the number of spatial cells, and assume
$v_{-k}=-v_k$ and $\omega_{-k}=\omega_k$.
We impose reflection at the endpoints and use reflected ghost values,
\[
g_{j,k}^n=g_{j,-k}^n\quad(j=0,J),\qquad
g_{-1,k}^n=g_{1,-k}^n,\qquad
g_{J+1,k}^n=g_{J-1,-k}^n.
\]
\begin{resultrevision}
Endpoint symmetry is imposed before the kinetic update, on the auxiliary variable $\widetilde g^{n+1}$, and on the provisional microscopic variable before checking positivity.
\end{resultrevision}
The diffusion and drift boundary fluxes are set to zero.
Thus the complete numerical density flux vanishes at both endpoints,
and the density update conserves mass.
Homogeneous Neumann conditions are used for the chemical variables.

For the correction, we assume $\varepsilon>0$, omit the superscript $n+1$,
and denote the provisional solution of
\eqref{eq:modified-main-scheme} by $\widehat\rho$ and $\widehat g$,
with $\langle\widehat g_{j,\cdot}\rangle_h=0$.
Define
\[
\widehat f_{j+\frac12,k}
=\widehat\rho_{j+\frac12}\psi_0(v_k)
+\frac{\varepsilon}{2}
 \bigl(\widehat g_{j,k}+\widehat g_{j+1,k}\bigr).
\]
\begin{resultrevision}
We first check the provisional density and the reconstructed distribution without modifying either variable. If both are nonnegative at every node, we return the provisional pair unchanged and skip all subsequent correction steps. Otherwise, we apply the correction below. Nonfinite inputs are rejected; no tolerance permits negative values to bypass this check.
\end{resultrevision}

For $\widehat\rho_{j+\frac12}>0$, set
\[
\theta_{j+\frac12}
=\min\left\{1,\,
\min_{\widehat f_{j+\frac12,k}<0}
\frac{\widehat\rho_{j+\frac12}\psi_0(v_k)}
{\widehat\rho_{j+\frac12}\psi_0(v_k)-\widehat f_{j+\frac12,k}}
\right\},
\]
where the inner minimum is omitted if all values are nonnegative, and
\[
\rho_{j+\frac12}^{*}=\widehat\rho_{j+\frac12},
\qquad
\bar g_{j+\frac12,k}^{*}
=\frac{\theta_{j+\frac12}}{2}
 \bigl(\widehat g_{j,k}+\widehat g_{j+1,k}\bigr).
\]
For $\widehat\rho_{j+\frac12}\le0$, set
$\rho_{j+\frac12}^{*}=\bar g_{j+\frac12,k}^{*}=0$.
Then $\rho_{j+\frac12}^{*}\psi_0(v_k)
+\varepsilon\bar g_{j+\frac12,k}^{*}\ge0$ and
$\langle\bar g_{j+\frac12,\cdot}^{*}\rangle_h=0$.
Assuming $\Delta x\sum_j\widehat\rho_{j+\frac12}>0$, restore the
input mass by
\[
\alpha=\frac{\sum_j\widehat\rho_{j+\frac12}}
                 {\sum_j\rho_{j+\frac12}^{*}},
\qquad
\rho_{j+\frac12}=\alpha\rho_{j+\frac12}^{*},
\qquad
\bar g_{j+\frac12,k}=\alpha\bar g_{j+\frac12,k}^{*}.
\]

We return the microscopic variable to interfaces by
\[
\begin{aligned}
\widehat h_{j,k}
 &=\tfrac12\bigl(\bar g_{j-\frac12,k}+\bar g_{j+\frac12,k}\bigr),
 &&1\le j\le J-1,\\
\widehat h_{0,k}
 &=\tfrac12\bigl(\bar g_{\frac12,k}+\bar g_{\frac12,-k}\bigr),&
\widehat h_{J,k}
 &=\tfrac12\bigl(\bar g_{J-\frac12,k}+\bar g_{J-\frac12,-k}\bigr),
\end{aligned}
\]
and set $h_j=(I-\Pi_h)\widehat h_j$.
Since this reconstruction need not preserve positivity, we check
\[
F_{j+\frac12,k}
=\rho_{j+\frac12}\psi_0(v_k)
+\frac{\varepsilon}{2}(h_{j,k}+h_{j+1,k}).
\]
If all values are nonnegative, we set $g_j=h_j$.
Otherwise, define
\[
\begin{gathered}
q_j=\frac{\min(\rho_{j-\frac12},\rho_{j+\frac12})}{|V|_h}
\quad(1\le j\le J-1),\qquad
q_0=\frac{\rho_{\frac12}}{|V|_h},\qquad
q_J=\frac{\rho_{J-\frac12}}{|V|_h},\\
m_j=\min_k h_{j,k},\qquad
\eta_j=
\begin{cases}
\min\{1,q_j/(-\varepsilon m_j)\},&m_j<0,\\
1,&m_j\ge0,
\end{cases}
\qquad g_{j,k}=\eta_j h_{j,k}.
\end{gathered}
\]
In this case, $\varepsilon g_{j,k}\ge-q_j$, and hence
\[
f_{j+\frac12,k}
=\rho_{j+\frac12}\psi_0(v_k)
+\frac{\varepsilon}{2}(g_{j,k}+g_{j+1,k})
\ge\frac{\rho_{j+\frac12}}{|V|_h}
-\frac{q_j+q_{j+1}}{2}\ge0.
\]
The interface scaling preserves the zero velocity mean and endpoint
reflection and leaves the density unchanged.
Together with the conservative density update, this gives
\[
\Delta x\sum_j\rho_{j+\frac12}
=\Delta x\sum_j\widehat\rho_{j+\frac12}
=\Delta x\sum_j\rho_{j+\frac12}^{\,n}.
\]
At fixed grids, this additional interface scaling is inactive for
sufficiently small $\varepsilon$ if $\rho$ is bounded below by a positive
constant independent of $\varepsilon$ and $h$ remains bounded.

\bibliographystyle{abbrv}
\bibliography{bib.bib}


\end{document}
